\documentclass{report}
\usepackage{amsmath,amssymb}
\setlength{\parindent}{0mm}
\setlength{\parskip}{1em}
\begin{document}
\begin{center}
\rule{6in}{1pt} \
{\large Klaudius Scheufele \\
{\bf Multi-Secant quasi-Newton Variants for Parallel Fluid-Structure Simulations -- and Other Multi-Physics Applications}}

Institute for Parallel and Distributed Systems \\ University of Stuttgart \\ Universit�tsstra�e 38 \\ D-70569 Stuttgart
\\
{\tt klaudius.scheufele@ipvs.uni-stuttgart.de}\\
Miriam Mehl\end{center}

Multi-secant quasi-Newton methods based on secant information produced
with no overhead throughout subsequent solver
iterations have been shown to be particularly suited to solve non-linear
fixed-point equations that arise from
partitioned multi-physics simulations where the exact Jacobian is
inaccessible. In all these methods, the multi-secant
equation for the approximate (inverse) Jacobian is enhanced by a norm
minimization condition. It is well-known, that
fluid-structure simulations, e.\,g., typically require the use of
secant-information from previous time steps. The number
of these time-steps highly depends on the application, its parameters,
the used solvers, and the mesh resolution. Using
to few leads to a relatively high number of iterations, using too many
not only to a computational overhead but also to
an increase of the number of iterations, as well. Determining the optimal
number requires a costly try-and-error process,
which can be avoided using a modified method that considers the difference of the
current (inverse) Jacobian and the one of the previous time step in the
norm-minimization. Thus, previous time step
information is taken into account in an implicit and automatized way
without magic parameters. We show numerical results
for fluid-structure interactions for both methods proving the robustness
and numerical efficiency of the second variant.
In addition, we propose a new algorithm eliminating the drawback of
having to store full interface Jacobian approximations
for the Jacobian difference norm minimization. This results in a highly
efficient, parallelizable, and robust iterative
solver applicable for surface coupling in many types of multi-physics simulations.


\end{document}
